%% kapa: exam=1 round=31 year=2020 period=1 form=A subject=economics subject_ko=경제학원론 seq=14 number=54
%% source: Q-Net 감정평가사 기출문제 (cst00309) artlSeq=5207219 file=extracted/1008/5207219/A형/A형/2020년도 제31회 감정평가사 1차 시험_1교시_A형.pdf
%% conversion: kordoc 4.24.51 -> tools/gen1.py (auto); review=proofread-sonnet; flags=none
\begin{question}{54}
甲의 소득은 24이고, $X$재와 $Y$재만 소비한다. 甲은 두 재화의 가격이 $P_{X}=4,P_{Y}=2$일 때 A( $x=5,y=1$)를 선택했고, 두 재화의 가격이 $P_{X}=3,P_{Y}=3$으로 변화함에 따라 B( $x=2,y=6$)를 선택했다. 甲의 선택에 관한 설명으로 옳은 것을 모두 고른 것은? (단, $x$는 $X$재 소비량, $y$는 $Y$재 소비량)
\begin{framedtext}
ㄱ. 甲은 가격 변화 전 B를 선택할 수 있었음에도 불구하고 A를 선택했다. \par
ㄴ. 甲은 가격 변화 후 A를 선택할 수 없었다. \par
ㄷ. 甲의 선택은 현시선호 약공리를 만족하지 못한다. \par
ㄹ. 甲은 주어진 예산제약 하에서 효용을 극대화하는 소비를 하고 있다. \par
\end{framedtext}
\begin{choices}
  \item ㄱ, ㄴ
  \item ㄱ, ㄷ
  \item ㄴ, ㄷ
  \item ㄴ, ㄹ
  \item ㄷ, ㄹ
\end{choices}
\end{question}
\answer{2}
